% ===========================================================================
%  sviluppatore/04-motore/01-barra.tex
% ===========================================================================
\section{Che cosa succede a ogni barra}\label{sez:barra}

Il ciclo principale sta in \file{backtest/engine.py}. Per ogni data del
calendario, dopo il periodo di riscaldamento, fa sei cose nello stesso ordine.

\figuraTikz{pipeline-motore}
  {I sei passi del ciclo. La divisione fra la riga superiore e quella
   inferiore non è grafica: è la divisione fra decidere ed eseguire.}
  {pipeline-motore}

\begin{enumerate}
  \item \textbf{Indicatori.} Non si calcolano dentro la strategia: si
        pre-calcolano una volta per simbolo e si passano già pronti. Che cosa
        serve lo dichiara ogni strategia nel proprio descrittore, e il
        registro ne fa l'unione.
  \item \textbf{Strategie.} Ciascuna riceve gli indicatori e il portafoglio e
        restituisce una lista di \simbolo{Signal}: che cosa, in che
        direzione, con quale convinzione. Non quanto.
  \item \textbf{Allocatore.} Dai segnali ai pesi di portafoglio. È qui che
        si decide quanto, e il come è configurabile
        (\vediSez{allocatore}).
  \item \textbf{Rischio.} Tetti su peso per posizione, esposizione lorda per
        classe di attività, leva complessiva, volatilità obiettivo.
  \item \textbf{Esecuzione.} Prezzo di apertura, slippage, costi.
  \item \textbf{Portafoglio.} Posizioni aggiornate, equity segnata a mercato,
        storia registrata.
\end{enumerate}

\subsection{Prima di tutto: i trigger intra-barra}

C'è un passo che precede la decisione, e non compare nella figura perché non
è parte del ciclo decisionale: il controllo degli \emph{stop} e dei
\emph{take profit} sulle posizioni già aperte.

\begin{codiceInterno}{python}{Il controllo che viene prima di tutto}{backtest/engine.py:745}
            # ---- INTRA-BAR TRIGGER CHECK ----
            # PRIMA della decision/execution phase, controlla se i SL/TP delle
            # posizioni aperte sono stati toccati nella bar t (high/low).
            trigger_hits = exit_tracker.check_bar(t, data)
\end{codiceInterno}

\noindent
Se un livello è stato toccato, la posizione si chiude subito e il simbolo
finisce in una lista nera fino al prossimo ribilanciamento vero. La filosofia
dichiarata nel codice è che «ribilanciare è quando rivaluti la tesi»: uno
stop toccato non aspetta il prossimo ribilanciamento, ma nemmeno autorizza a
riaprire la stessa posizione nella stessa barra.

\begin{attenzione}{La lista nera non blocca le chiusure}
Un simbolo in lista nera non può essere \emph{aperto}, ma un target a zero
passa sempre. Invertire questa asimmetria --- cioè bloccare anche le chiusure
--- significherebbe tenere aperta una posizione che il sistema ha deciso di
chiudere, ed è il tipo di errore che si scopre su un conto reale.
\end{attenzione}
